% ============================================================
%  MODULE VII — Risk Management
% ============================================================

\chapter{Risk Management}

\begin{notebox}
\textbf{Risk management is what keeps you in the game.} Every professional trader's primary
goal is capital preservation — the ability to keep playing. Edge compounds over time only
if you survive the drawdowns. The best entry signal in the world is worthless without a
disciplined risk framework.
\end{notebox}

% ────────────────────────────────────────────────────────────
\section{Drawdown}

\begin{keyterm}{Drawdown (DD)}
The peak-to-trough decline in portfolio value, expressed as a percentage.
\[
  DD = \frac{P_{\text{peak}} - P_{\text{trough}}}{P_{\text{peak}}} \times 100\%
\]
A 50\% drawdown requires a \textbf{100\% return} just to get back to breakeven.
\end{keyterm}

\begin{center}
\begin{tabular}{>{\bfseries}p{3.5cm} p{4cm} p{5cm}}
\toprule
Loss & Recovery Required & Implied Effect \\
\midrule
10\% & 11.1\% & Manageable \\
20\% & 25.0\% & Uncomfortable \\
30\% & 42.9\% & Serious \\
40\% & 66.7\% & Very serious \\
50\% & 100.0\% & Account in crisis \\
75\% & 300.0\% & Near impossible \\
\bottomrule
\end{tabular}
\end{center}

This is why limiting drawdowns is more important than maximising returns.

% ────────────────────────────────────────────────────────────
\section{Portfolio-Level Risk Rules}

\subsection{Daily / Weekly Loss Limits}
Set a \textbf{circuit breaker}: if you lose $X$\% in a day (e.g., 2\%), stop trading.
Come back tomorrow. This prevents emotional spiralling from turning a bad day into a catastrophe.

\subsection{Maximum Open Risk}
At any given time, total open risk across all positions should be capped. Example:
no more than 5--6\% of account at risk simultaneously.

\subsection{Correlation Risk}
Holding 10 positions in the same sector = effectively one concentrated bet.
True diversification requires \textbf{uncorrelated} positions.

% ────────────────────────────────────────────────────────────
\section{Value at Risk (VaR)}

\begin{keyterm}{Value at Risk (VaR)}
The maximum expected loss over a given time horizon at a specified confidence level.
``95\% 1-day VaR of \$10{,}000'' means: there is a 5\% chance of losing more than
\$10{,}000 in one day.
\end{keyterm}

\begin{warningbox}
VaR fails precisely when you need it most — in fat-tail events and market crises,
correlations spike to 1 and losses exceed VaR estimates dramatically.
The 2008 financial crisis exposed VaR's limitations. Use it as one tool, not the only tool.
\end{warningbox}

% ────────────────────────────────────────────────────────────
\section{The Kelly Criterion}

Optimal position sizing formula to maximise long-run portfolio growth.

\begin{formulabox}
\[
  f^* = \frac{p}{a} - \frac{q}{b}
\]
where:
\begin{itemize}[noitemsep]
  \item $f^*$ = fraction of capital to risk
  \item $p$ = probability of winning
  \item $q$ = probability of losing $(= 1 - p)$
  \item $b$ = net profit per unit risked (win/loss ratio)
  \item $a$ = loss per unit risked (= 1 in standard form)
\end{itemize}
\end{formulabox}

\begin{examplebox}
Win rate = 55\%, Win/Loss ratio = 1.5.
\[
  f^* = \frac{0.55}{1} - \frac{0.45}{1.5} = 0.55 - 0.30 = 0.25 = 25\%
\]
Full Kelly says bet 25\% of capital. In practice, use \textbf{Half Kelly} (12.5\%) to reduce
volatility while capturing most of the growth benefit.
\end{examplebox}

\begin{notebox}
Kelly sizing assumes you know your true edge precisely — rarely the case in markets.
Over-betting Kelly leads to catastrophic drawdowns. Most professional traders
use a fraction of Kelly (25--50\%).
\end{notebox}

% ────────────────────────────────────────────────────────────
\section{Sharpe Ratio and Performance Metrics}

\begin{formulabox}
\textbf{Sharpe Ratio:}
\[
  S = \frac{R_p - R_f}{\sigma_p}
\]
where $R_p$ = portfolio return, $R_f$ = risk-free rate, $\sigma_p$ = standard deviation of returns.
\end{formulabox}

A Sharpe Ratio of 1.0 is decent; 2.0 is very good; 3.0+ is exceptional (most quant funds target this).

\begin{center}
\begin{tabular}{>{\bfseries}p{4cm} p{8.5cm}}
\toprule
Metric & What It Measures \\
\midrule
Sharpe Ratio & Return per unit of total volatility \\
Sortino Ratio & Like Sharpe but only penalises \emph{downside} volatility \\
Calmar Ratio & Annualised return $\div$ Max Drawdown \\
Win Rate & \% of trades that are profitable \\
Profit Factor & Gross profit $\div$ Gross loss (want $>1.5$) \\
Expectancy & Average \$ made per trade (positive = edge) \\
\bottomrule
\end{tabular}
\end{center}

% ────────────────────────────────────────────────────────────
\subsection{Benchmark-Relative Performance Metrics}

The metrics above measure performance in isolation. A separate, equally important
question is: \emph{how did you do compared to the market?}
A portfolio can fall 3\% on a day the market falls 5\% — that is a loss in absolute
terms, but an outperformance in relative terms. The following metrics formalise this idea.

\begin{keyterm}{Excess Return (Active Return)}
The simplest benchmark-relative metric. It answers: by how many percentage points did
you beat or lag the market?
\[
  r_{\text{excess}} = r_{\text{you}} - r_{\text{market}}
\]
\textbf{Example:} You are down $-3\%$ and the market is down $-5\%$:
\[
  -3\% - (-5\%) = +2\%
\]
You lost money in absolute terms, but outperformed the benchmark by 2 percentage points.
\end{keyterm}

\begin{keyterm}{Relative Return Ratio}
A normalised ratio version, useful when returns are large and the difference alone
can be misleading.
\[
  r_{\text{relative}} = \frac{1 + r_{\text{you}}}{1 + r_{\text{market}}} - 1
\]
\textbf{Example:} You return $-3\%$ against a market of $-5\%$:
\[
  \frac{0.97}{0.95} - 1 \approx +2.11\%
\]
Your portfolio performed about $2.1\%$ better at the level of compounded returns.
\end{keyterm}

\begin{keyterm}{Beta-Adjusted Abnormal Return}
If your portfolio tends to move more (or less) than the market due to its intrinsic
volatility exposure, comparing raw returns is unfair. \textbf{Beta} ($\beta$) measures
your sensitivity: $\beta = 2$ means your portfolio typically moves twice as much as
the index.
\[
  r_{\text{abnormal}} = r_{\text{you}} - \beta \cdot r_{\text{market}}
\]
\textbf{Example:} $\beta = 1.5$, market returns $-4\%$, you return $-4\%$:
\[
  -4\% - 1.5 \times (-4\%) = -4\% + 6\% = +2\%
\]
You were expected to fall 6\% given your beta; falling only 4\% is a 2\% positive
abnormal return. This is the foundation of \textbf{alpha}.
\end{keyterm}

\begin{keyterm}{Capture Ratio (Upside / Downside)}
Measures what fraction of the market's moves — up or down — your portfolio captures.
\[
  \text{Downside Capture} = \frac{r_{\text{you, down markets}}}{r_{\text{market, down markets}}}
\]
\[
  \text{Upside Capture} = \frac{r_{\text{you, up markets}}}{r_{\text{market, up markets}}}
\]
A capture ratio below 1.0 on the downside is good (you lose less than the market
during drawdowns). A ratio above 1.0 on the upside is good (you gain more than the
market in rallies).

\textbf{Example — Downside:} Market falls $-10\%$, you fall $-6\%$:
\[
  \frac{-6\%}{-10\%} = 0.6
\]
You captured only 60\% of the market's downside — good.

\textbf{Example — Downside (bad):} Market falls $-10\%$, you fall $-15\%$:
\[
  \frac{-15\%}{-10\%} = 1.5
\]
You lost 50\% worse than the market.
\end{keyterm}

\begin{notebox}
\textbf{Alpha} is the umbrella term for risk-adjusted outperformance of a benchmark.
In its simplest form, alpha is the intercept in a regression of your returns against
market returns:
\[
  r_{\text{you}} = \alpha + \beta \cdot r_{\text{market}} + \varepsilon
\]
Positive alpha means you generate excess return that cannot be explained by
simply holding a leveraged position in the market.
\end{notebox}

\begin{center}
\begin{tabular}{>{\bfseries}p{4.5cm} p{7.5cm}}
\toprule
Metric & Best used when \ldots \\
\midrule
Excess return & Quick comparison; same volatility profile as benchmark \\
Relative return ratio & Large return magnitudes; compounding matters \\
Beta-adjusted abnormal return & Portfolio has different volatility exposure than index \\
Downside capture ratio & Evaluating defensive quality of strategy in down markets \\
Alpha & Full risk-adjusted, regression-based performance attribution \\
\bottomrule
\end{tabular}
\end{center}

% ────────────────────────────────────────────────────────────
\section{Hedging}

Reducing existing risk by taking an offsetting position.

\begin{itemize}
  \item \textbf{Long stock + Buy Put option} — limits downside loss on the stock position.
  \item \textbf{Long stock + Short futures} — locks in a price; used by producers and large funds.
  \item \textbf{Long equities + Long VIX calls} — hedge against volatility spikes in crises.
  \item \textbf{Pair trade} — long one stock, short a correlated stock in same sector.
        Removes market beta; profits from relative performance.
\end{itemize}

Hedges have a cost (premium, drag, opportunity cost). Perfect hedges eliminate profit
as well as loss. In practice, \textbf{partial hedges} are used.

% ────────────────────────────────────────────────────────────
\section{The Trader's Ruin Problem}

Even with a positive edge, a trader can go bankrupt through bad luck and poor sizing.

\begin{formulabox}
\textbf{Probability of Ruin (simplified):}
\[
  P(\text{ruin}) = \left(\frac{q}{p}\right)^N
\]
where $N$ = number of units of risk in the account, $p$ = win probability, $q$ = loss probability.
\end{formulabox}

The message: \textbf{adequate capitalisation} combined with \textbf{small risk per trade}
dramatically reduces the probability of ruin, even in a random walk with slight positive edge.

\begin{warningbox}
Never risk money you cannot afford to lose. Do not trade with rent money, emergency funds,
or borrowed capital. The psychological pressure of trading money you \emph{need} destroys
decision-making and guarantees losses.
\end{warningbox}
